% GENERATED FILE. Edit canonical lessons, then run npm run generate:books.
% canonical-source-hash: fnv1a64:e131c682791a22c3
% canonical-lessons: GU-C249-mobail, GU-C249-nambar, GU-C249-sandesho, GU-C249-avaj, GU-C249-melavvu

\chapter{Calls and messages}
\label{ch:gu-calls-and-messages}
\hlchaptermodality{249}

\begin{chapteropening}
\textbf{By the end of this chapter:} \emph{I can talk about my mobile phone, numbers, messages and voices.}

Complete the last lesson of chapter 249: I can talk about my mobile phone, numbers, messages and voices.
\end{chapteropening}

\section[\texorpdfstring{\gu{મોબાઈલ} (mobāīl)}{મોબાઈલ (mobāīl)}]{\gu{મોબાઈલ} (mobāīl) — mobile phone}

\label{lesson:GU-C249-mobail}



\begin{quote}
Before the new one: say the Gujarati for to hand over, then the Gujarati for to note down.
\end{quote}

\subsection*{You'll want to know: \gu{મોબાઈલ}}
\textbf{\gu{મોબાઈલ}} — \emph{mobāīl} — \textquotedblleft{}mobile phone\textquotedblright{}.

\begin{quote}
\textbf{\gu{મારો} \gu{મોબાઈલ} \gu{ખિસ્સામાં} \gu{છે}.} — \emph{māro mobāīl khissāmā̃ che} — \textquotedblleft{}My mobile phone is in my pocket.\textquotedblright{}
\end{quote}

\begin{cousinweb}
\textbf{\gu{મોબાઈલ}} (\emph{mobāīl}) is a loan from English \emph{mobile}, short for a mobile phone.
\end{cousinweb}

\begin{grammarlens}[title={What you've built}]
One more word for everyday life at A2.
\end{grammarlens}

\subsection*{Your turn}
\begin{itemize}
\raggedright
  \item \emph{Say it:} \emph{mobāīl}
  \item \emph{Say it:} \emph{mobāīl}, once more
  \item \emph{Say it:} \emph{māro mobāīl khissāmā̃ che} — My mobile phone is in my pocket.
  \item \emph{Recall:} say the Gujarati for to hand over, then the Gujarati for to note down, then say \emph{mobāīl} again
\end{itemize}

\begin{tcolorbox}[breakable,colback=teal!4,colframe=teal!35!black,title={Before you move on}]
What does \gu{મોબાઈલ} mean? (\textquotedblleft{}Mobile phone\textquotedblright{}.) Say it once more.
\end{tcolorbox}
\section[\texorpdfstring{\gu{નંબર} (nambar)}{નંબર (nambar)}]{\gu{નંબર} (nambar) — number}

\label{lesson:GU-C249-nambar}



\begin{quote}
Before the new one: say the Gujarati for to note down, then the Gujarati for mobile phone.
\end{quote}

\subsection*{You'll want to know: \gu{નંબર}}
\textbf{\gu{નંબર}} — \emph{nambar} — \textquotedblleft{}number\textquotedblright{}.

\begin{quote}
\textbf{\gu{તમારો} \gu{ફોન} \gu{નંબર} \gu{શું} \gu{છે}?} — \emph{tamāro phon nambar śũ che?} — \textquotedblleft{}What is your phone number?\textquotedblright{}
\end{quote}

\begin{cousinweb}
\textbf{\gu{નંબર}} (\emph{nambar}) is a loan from English \emph{number}, used for phone, house and bus numbers.
\end{cousinweb}

\begin{grammarlens}[title={What you've built}]
One more word for everyday life at A2.
\end{grammarlens}

\subsection*{Your turn}
\begin{itemize}
\raggedright
  \item \emph{Say it:} \emph{nambar}
  \item \emph{Say it:} \emph{nambar}, once more
  \item \emph{Say it:} \emph{tamāro phon nambar śũ che?} — What is your phone number?
  \item \emph{Recall:} say the Gujarati for to note down, then the Gujarati for mobile phone, then say \emph{nambar} again
\end{itemize}

\begin{tcolorbox}[breakable,colback=teal!4,colframe=teal!35!black,title={Before you move on}]
What does \gu{નંબર} mean? (\textquotedblleft{}Number\textquotedblright{}.) Say it once more.
\end{tcolorbox}
\section[\texorpdfstring{\gu{સંદેશો} (sandeśo)}{સંદેશો (sandeśo)}]{\gu{સંદેશો} (sandeśo) — message}

\label{lesson:GU-C249-sandesho}



\begin{quote}
Before the new one: say the Gujarati for mobile phone, then the Gujarati for number.
\end{quote}

\subsection*{You'll want to know: \gu{સંદેશો}}
\textbf{\gu{સંદેશો}} — \emph{sandeśo} — \textquotedblleft{}message\textquotedblright{}.

\begin{quote}
\textbf{\gu{મારો} \gu{ભાઈ} \gu{રોજ} \gu{સંદેશો} \gu{મોકલે} \gu{છે}.} — \emph{māro bhāī roj sandeśo mokle che} — \textquotedblleft{}My brother sends a message every day.\textquotedblright{}
\end{quote}

A masculine noun.

\begin{cousinweb}
\textbf{\gu{સંદેશો}} (\emph{sandeśo}) comes from Sanskrit \emph{saṁdēśa}, \textquotedblleft{}message, information\textquotedblright{} (CDIAL 12904). It keeps the Sanskrit shape with the Gujarati masculine ending \emph{-o}, a semi-tatsama.
\end{cousinweb}

\begin{grammarlens}[title={What you've built}]
One more word for everyday life at A2.
\end{grammarlens}

\subsection*{Your turn}
\begin{itemize}
\raggedright
  \item \emph{Say it:} \emph{sandeśo}
  \item \emph{Say it:} \emph{sandeśo}, once more
  \item \emph{Say it:} \emph{māro bhāī roj sandeśo mokle che} — My brother sends a message every day.
  \item \emph{Recall:} say the Gujarati for mobile phone, then the Gujarati for number, then say \emph{sandeśo} again
\end{itemize}

\begin{tcolorbox}[breakable,colback=teal!4,colframe=teal!35!black,title={Before you move on}]
What does \gu{સંદેશો} mean? (\textquotedblleft{}Message\textquotedblright{}.) Say it once more.
\end{tcolorbox}
\section[\texorpdfstring{\gu{અવાજ} (avāj)}{અવાજ (avāj)}]{\gu{અવાજ} (avāj) — voice, sound}

\label{lesson:GU-C249-avaj}



\begin{quote}
Before the new one: say the Gujarati for number, then the Gujarati for message.
\end{quote}

\subsection*{You'll want to know: \gu{અવાજ}}
\textbf{\gu{અવાજ}} — \emph{avāj} — \textquotedblleft{}voice, sound\textquotedblright{}.

\begin{quote}
\textbf{\gu{ફોનમાં} \gu{મારી} \gu{માતાનો} \gu{અવાજ} \gu{આવે} \gu{છે}.} — \emph{phonmā̃ mārī mātāno avāj āve che} — \textquotedblleft{}My mother's voice comes through the phone.\textquotedblright{}
\end{quote}

\begin{cousinweb}
\textbf{\gu{અવાજ}} (\emph{avāj}) is a loan, usually traced to Persian \emph{āvāz}, \textquotedblleft{}voice, sound\textquotedblright{}. CDIAL has no entry for it.
\end{cousinweb}

\begin{grammarlens}[title={What you've built}]
One more word for everyday life at A2.
\end{grammarlens}

\subsection*{Your turn}
\begin{itemize}
\raggedright
  \item \emph{Say it:} \emph{avāj}
  \item \emph{Say it:} \emph{avāj}, once more
  \item \emph{Say it:} \emph{phonmā̃ mārī mātāno avāj āve che} — My mother's voice comes through the phone.
  \item \emph{Recall:} say the Gujarati for number, then the Gujarati for message, then say \emph{avāj} again
\end{itemize}

\begin{tcolorbox}[breakable,colback=teal!4,colframe=teal!35!black,title={Before you move on}]
What does \gu{અવાજ} mean? (\textquotedblleft{}Voice, sound\textquotedblright{}.) Say it once more.
\end{tcolorbox}
\section[\texorpdfstring{\gu{મેળવવું} (meḷavvũ)}{મેળવવું (meḷavvũ)}]{\gu{મેળવવું} (meḷavvũ) — to obtain, to get}

\label{lesson:GU-C249-melavvu}



\begin{quote}
Before the new one: say the Gujarati for message, then the Gujarati for voice.
\end{quote}

\subsection*{You'll want to know: \gu{મેળવવું}}
\textbf{\gu{મેળવવું}} — \emph{meḷavvũ} — \textquotedblleft{}to obtain, to get\textquotedblright{}.

\begin{quote}
\textbf{\gu{હું} \gu{મિત્ર} \gu{પાસેથી} \gu{નંબર} \gu{મેળવું} \gu{છું}.} — \emph{hũ mitra pāsethī nambar meḷavũ chũ} — \textquotedblleft{}I get the number from a friend.\textquotedblright{}
\end{quote}

\begin{cousinweb}
\textbf{\gu{મેળવવું}} (\emph{meḷavvũ}), to bring something to oneself, is usually explained as the causative of the inherited \emph{maḷvũ}, \textquotedblleft{}to meet, to be got\textquotedblright{}, with the vowel raised to \emph{e}. CDIAL derives \emph{maḷvũ} from Sanskrit \emph{miláti}, \textquotedblleft{}meets\textquotedblright{} (CDIAL 10133), but lists no Gujarati form under the Sanskrit causative \emph{mēlayati} (CDIAL 10332).
\end{cousinweb}

\begin{grammarlens}[title={What you've built}]
One more word for everyday life at A2.
\end{grammarlens}

\subsection*{Your turn}
\begin{itemize}
\raggedright
  \item \emph{Say it:} \emph{meḷavvũ}
  \item \emph{Say it:} \emph{meḷavvũ}, once more
  \item \emph{Say it:} \emph{hũ mitra pāsethī nambar meḷavũ chũ} — I get the number from a friend.
  \item \emph{Recall:} say the Gujarati for message, then the Gujarati for voice, then say \emph{meḷavvũ} again
\end{itemize}

\begin{tcolorbox}[breakable,colback=teal!4,colframe=teal!35!black,title={Before you move on}]
What does \gu{મેળવવું} mean? (\textquotedblleft{}To obtain, to get\textquotedblright{}.) Say it once more.
\end{tcolorbox}
